\title{FlashAttention-3: Fast and Accurate Attention with Asynchrony and Low-precision}

\begin{abstract}
The attention mechanism, fundamental to the widespread Transformer architecture, remains a major computational bottleneck in scaling large language models and handling extended contexts. While \textsc{FlashAttention} previously accelerated attention by reducing memory accesses on GPUs, it does not fully utilize enhancements introduced by the latest hardware; notably, \textsc{FlashAttention-2} manages only 35\% peak throughput on the H100 GPU. We introduce three strategies to further enhance attention computation on Hopper GPUs: leveraging asynchrony between Tensor Cores and TMA to (1) concurrently perform computation and memory transfers via warp specialization, (2) fuse block-wise matrix multiplication with softmax operations for increased parallelism, and (3) employ block quantization alongside incoherent processing to harness hardware-enabled FP8 low-precision support. Our proposed approach, \textsc{FlashAttention-3}, achieves 1.5-2.0$\times$ faster execution on H100 GPUs, attaining up to 740 TFLOPs/s with FP16 (corresponding to 75\% hardware utilization) and nearly 1.2 PFLOPs/s with FP8 precision. Additionally, we show that the FP8 implementation of \textsc{FlashAttention-3} reduces numerical error by a factor of 2.6 compared to a reference FP8 attention baseline.
\end{abstract}

\section{Introduction}
\label{sec:intro}

Within the Transformer framework~\citep{vaswani2017attention}, the central performance limitation arises from the attention mechanism, specifically due to the quadratic computational complexity involved in determining self-attention scores based on pairs of queries and keys as sequence length increases. Addressing this bottleneck to accommodate much longer contexts could facilitate advancements in several areas: it enables reasoning and modeling across extended textual passages~\citep{guo2021longt5,shaham2022scrolls,peng2023yarn} and managing multiple files within extensive code repositories~\citep{roziere2023code, li2023starcoder}; it opens up further potential for new data modalities, such as processing high-resolution visual data~\citep{chen2022scaling}, audio streams~\citep{gulati2020conformer}, and video content~\citep{ho2022video}; and it also introduces possibilities for expanded applications, like maintaining coherent user histories in interactive systems~\citep{sun2019bert4rec} or supporting agents with long-term operational planning~\citep{yao2022react}. Consequently, there is a burgeoning research effort to accelerate attention computations for extended input contexts. This includes exploration of approximate attention methods~\citep{katharopoulos2020transformers,choromanski2020rethinking, tay2020efficient}, advancements in computational efficiency through software~\citep{rabe2021self,dao2022flashattention,kwon2023efficient}, as well as investigations into entirely new neural architectures~\citep{peng2023rwkv,sun2023retentive,gu2023mamba}.

This paper extends the foundation established by \citet{dao2022flashattention}, who developed precise attention algorithms by embedding an awareness of GPU execution patterns and hardware architecture directly into their algorithmic framework. In their seminal work, Dao et al. presented \textsc{FlashAttention}, a groundbreaking approach that utilizes tiling to enable efficient parallel execution of attention, thereby circumventing costly transfers to and from global memory by merging all attention computations into a single GPU kernel invocation.

Subsequently, \citet{dao2023flashattention2} advanced this strategy with \textsc{FlashAttention-2}, revising the algorithm to also exploit parallelism along the sequence length and to process blocks from the key and value matrices within the inner loop of the forward pass. This adjustment aimed to maximize GPU resource usage and evenly distribute computational loads. Despite these innovations, our findings indicate that \textsc{FlashAttention-2} still falls short of harnessing the full computational capacity of contemporary GPUs—yielding approximately 35\% utilization compared to 80–90\% observed with finely tuned matrix-multiplication (GEMM) routines on devices such as the Hopper H100. This shortfall may partly stem from implementation choices, including the use of instructions tailored to earlier GPU generations (e.g., Ampere) instead of those designed for Hopper Tensor Cores. Prior works like ThunkerKitten~\citep{spector2024thunder} and cuDNN 9~\citep{cudnn9} have demonstrated that leveraging Hopper-specific instructions, together with a tile-centric programming paradigm, can yield substantial improvements both in computation speed and implementation simplicity for attention mechanisms.

At a deeper level, the algorithm underlying \textsc{FlashAttention-2} operates within a simplified, fully synchronous computation framework, deliberately omitting explicit mechanisms for leveraging asynchrony or reduced precision during execution. In practice, asynchrony arises from hardware-level differentiation that expedites the most critical routines in machine learning pipelines: dedicated units for matrix multiplication (Tensor Cores) and specialized modules for memory accesses (Tensor Memory Accelerator -- TMA), both working independently from the generic CUDA cores that handle logic and numerical tasks. Meanwhile, adoption of low precision formats—such as FP8 in Hopper and FP4 in Blackwell, in succession to the earlier employment of FP16 (first introduced with Pascal in 2017) and BF16 (with Ampere in 2020)—has demonstrably increased computational throughput (by factors of two or four) for fixed silicon area and power usage. The features and advantages specific to Hopper in terms of both asynchrony and precision are further surveyed in \cref{subsec:hardware}. 
The core engineering obstacle lies in revisiting the design of \textsc{FlashAttention-2} to effectively utilize these modern hardware capabilities: asynchrony introduces the necessity for interleaving operations like matrix multiplication and softmax (despite their data dependencies), while low-precision support demands methods to control quantization errors, which are particularly impactful in the processing of outlier features common to LLMs~\citep{dettmers2208llm, sun2024massive}.

In response, we introduce \textsc{FlashAttention-3}, presenting a synthesis of three novel conceptual advances aimed at maximizing efficiency on contemporary GPU platforms.\footnote{While our empirical analysis targets NVIDIA’s Hopper generation, the proposed methods extend straightforwardly to any GPU offering comprehensive asynchronous and reduced-precision functionality.}

\begin{enumerate}

\item \textbf{Producer-Consumer asynchrony:} We introduce a warp-specialized software pipelining approach that capitalizes on the decoupled execution of data transfers and Tensor Core computations. By assigning producer and consumer roles for data to distinct warps, our scheme enhances the algorithm’s effectiveness at masking memory and instruction latency.
\item \textbf{Hiding softmax under asynchronous block-wise GEMMs:} The low-throughput, non-GEMM tasks required for softmax—such as floating point multiply-add and exponential operations—are scheduled to run concurrently with asynchronous WGMMA-based GEMM computations. This necessitates a refactoring of the \textsc{FlashAttention-2} procedure to eliminate select sequential dependencies that typically exist between softmax and GEMM steps. For instance, in our algorithm’s two-stage variant, one warp computes softmax on a block of the score matrix, while the WGMMA instruction asynchronously prepares the subsequent block.
\item \textbf{Hardware-accelerated low-precision GEMM:} We revise the forward computation method to utilize the FP8 Tensor Cores for GEMM, which leads to an almost twofold increase in observed TFLOPs/s. Accommodating the divergent memory layout expectations for FP32 accumulators and FP8 inputs in WGMMA involves employing strategies such as block quantization and incoherent processing, thereby compensating for the precision loss incurred when transitioning to FP8.

To empirically assess the effectiveness of our approach, we evaluate \textsc{FlashAttention-3} on the H100 SXM5 GPU across diverse parameter settings. Our results demonstrate that (1) in FP16 precision, \textsc{FlashAttention-3} achieves a forward pass speedup of 1.5-2.0$\times$ over \textsc{FlashAttention-2}, reaching peak performance up to 740 TFLOPs/s, and achieves 1.5-1.75$\times$ acceleration in the backward pass; (2) when running in FP8, the method attains nearly 1.2 PFLOPs/s; and (3) for scenarios involving long sequence lengths, FP16 consistently exceeds, and FP8 closely matches\footnote{Specifically, with head dimension 64, \textsc{FlashAttention-3} in FP8 surpasses other methods, while for head dimensions 128 and 256 it performs comparably in the absence of causal masking but trails behind when causal masking is applied.} the performance of the leading attention implementation provided by NVIDIA's cuDNN. Additionally, we confirm that FP16 \textsc{FlashAttention-3} incurs no greater numerical error than \textsc{FlashAttention-2}, and in fact demonstrates improved accuracy relative to the baseline attention implementation by performing intermediate computations (such as softmax rescaling) in FP32. Notably, our FP8 implementation utilizing block quantization combined with incoherent processing exhibits a 2.6$\times$ accuracy improvement over the conventional attention method that uses per-tensor quantization when processing features with significant outliers.

\textsc{FlashAttention-3} is released as open source under a permissive license\footnote{\textsc{FlashAttention-3}
  is available at \texttt{https://github.com/Dao-AILab/flash-attention}}, with plans for integration into both the PyTorch and Hugging Face ecosystems to maximize accessibility and utility for researchers and practitioners.

\section{Background: Multi-Head Attention and GPU Characteristics}
\label{sec:background}

\subsection{Multi-Head Attention}
\label{subsec:multi_head_attn}

Consider $\mathbf{Q}, \mathbf{K}, \mathbf{V} \in \mathbb{R}^{N \times d}$, which denote the query, key, and value matrices for a given attention head. Here, $N$ represents the sequence length and $d$ indicates the dimensionality of each head. The output of the attention mechanism, denoted by $\mathbf{O}$, is obtained as follows:

\begin{equation*}
  \mathbf{S} = \alpha \mathbf{Q} \mathbf{K}^\top \in \mathbb{R}^{N \times N}, \quad \mathbf{P} = \mathrm{softmax}(\mathbf{S}) \in \mathbb{R}^{N \times N}, \quad \mathbf{O} = \mathbf{P}\mathbf{V} \in \mathbb{R}^{N \times d},
\end{equation*}

The $\mathrm{softmax}$ function operates along each row, and it is common to use a scaling coefficient of $\alpha = 1/\sqrt{d}$. To improve numerical stability when computing exponentials, we typically subtract $\mathrm{rowmax}(\mathbf{S})$ from $\mathbf{S}$. In the context of multi-head attention (MHA), separate query, key, and value projections are assigned to each head. These calculations are executed in parallel across all heads and input batches, resulting in the complete output tensor.

Let $\phi$ denote a scalar loss function, and use the shorthand $\mathbf{d}(-) = \partial \phi / \partial (-)$ for the gradient operator. Starting from the gradient of the output, $\mathbf{dO} \in \mathbb{R}^{N \times d}$, we derive the gradients $\mathbf{dQ}$, $\mathbf{dK}$, and $\mathbf{dV}$ by sequentially applying the chain rule, as shown below:

\begin{align*}
  \mathbf{dV} &= \mathbf{P}^\top \mathbf{dO} \in \mathbb{R}^{N \times d} \\
  \mathbf{dP} &= \mathbf{dO} \mathbf{V}^\top \in \mathbb{R}^{N \times N} \\
  \mathbf{dS} &= \mathrm{dsoftmax} (\mathbf{dP}) \in \mathbb{R}^{N \times N} \\
  \mathbf{dQ} &= \alpha \mathbf{dS} \mathbf{K} \in \mathbb{R}^{N \times d} \\
  \mathbf{dK} &= \alpha \mathbf{dS}^\top \mathbf{Q} \in \mathbb{R}^{N \times d},
\end{align*}

In these equations, we define the Jacobian of softmax as $\mathbf{d}s = (\mathrm{diag}(p) - p p^\top)\mathbf{d}p$, where $p = \mathrm{softmax}(s)$ is the softmax function taken over vector $s$. We interpret $\mathrm{dsoftmax}(\mathbf{dP})$ as the result of applying this operation independently to each row. Finally, this procedure allows the backward pass for MHA to be carried out in parallel over both heads and batch dimensions.

\subsection{GPU hardware characteristics and execution model}
\label{subsec:hardware}

We describe the aspects of the GPU's execution model relevant for \textsc{FlashAttention-3}, with a focus on the NVIDIA Hopper architecture as a concrete instantiation of this model.

\paragraph{Memory hierarchy:} The memory architecture of the GPU comprises several layers of data storage, structured such that bandwidth increases as capacity decreases (\cref{tab:gpu-hierarchy})\footnote{\citet{luo2024benchmarking} identifies that the shared memory bandwidth is 128 bytes per cycle per SM; this value is multiplied by 132 SMs and the boost clock speed of 1830 MHz.}. At the lowest level is global memory (GMEM), or HBM, which refers to the external DRAM shared by all streaming multiprocessors (SMs). Whenever data is fetched from GMEM, it is automatically stored in an intermediate on-chip L2 cache. Each SM is also equipped with a limited-size, multi-banked, programmer-controlled cache known as shared memory (SMEM). Finally, registers are housed within every SM, constituting the fastest but most limited level of the hierarchy.

\paragraph{Thread hierarchy:} In the GPU programming paradigm, execution units known as threads are structured into various logical groupings. This hierarchy ranges from the most granular to the most aggregate levels, starting with individual threads, followed by warps (each containing 32 threads), then warpgroups (which consist of 4 sequential warps), and next by threadblocks (also referred to as cooperative thread arrays, or CTAs). Further aggregation in Hopper includes threadblock clusters, with grids representing the highest level of this organizational structure.

The relationship between these hierarchies is highly integrated. Threads that belong to a single CTA are jointly scheduled on an identical SM, while CTAs within a particular cluster are concurrently executed on the same GPC. All threads inside a CTA have direct access to SMEM, but each individual thread can utilize up to 256 private registers (RMEM).

\paragraph{Asynchrony and warp-specialization:}

GPUs achieve high throughput by leveraging both parallelism and asynchronous execution, enabling them to mask delays associated with memory access and computation. In the Hopper architecture, the Tensor Memory Accelerator (TMA) serves as specialized hardware to facilitate asynchronous data transfers between GMEM and SMEM \cite[\S7.29]{cuda}. Moreover, Hopper introduces an asynchronous Tensor Core—accessed through the warpgroup-level WGMMA instruction \cite[\S9.7.14]{ptx}—which, in contrast to earlier architectures like Ampere, is capable of fetching its operands directly from shared memory.

The inclusion of hardware mechanisms to facilitate asynchrony enables the creation of warp-specialized kernels, in which different warps within a CTA are assigned either producer or consumer tasks, resulting in warps dedicated solely to data transfer or computation operations. This general approach enhances the compiler's capacity to produce more efficient instruction orderings \citep{warp-specialization-2011}. Furthermore, Hopper provides a feature for on-the-fly reallocation of registers among warpgroups through the \verb|setmaxnreg| directive \cite[\S9.7.17.1]{ptx}, permitting warps handling MMAs to access a larger fraction of RMEM compared to warps restricted to TMA operations, which typically require only one active thread.

\paragraph{Low-precision number formats:}
\label{sec:low-precision-gpu}
Contemporary GPUs come equipped with specialized units designed to boost the efficiency of low-precision arithmetic. For instance, Hopper’s WGMMA instruction can utilize FP8 Tensor Cores, achieving double the SM throughput relative to what is attainable with FP16 or BF16 formats.

Correctly utilizing FP8 WGMMA requires a clear grasp of the operand layout requirements. In the context of a GEMM operation that computes $A \times B^{\top}$, where $A$ is an $M\times K$ matrix and $B$ is an $N\times K$ matrix, the operand $A$ or $B$ is considered \emph{mn-major} if its elements are stored contiguously along the outer $M$ or $N$ dimension, respectively. Conversely, the operand is \emph{k-major} if its storage is contiguous along the inner $K$-dimension. For FP16 WGMMA, SMEM operands can be provided in either mn-major or k-major formats. In contrast, FP8 WGMMA only permits k-major formatted inputs. Additionally, challenges arise in scenarios such as attention, where chaining multiple GEMM operations within a single kernel is desirable; specifically, inconsistencies between FP32 accumulator layouts and FP8 operand layouts complicate the invocation of dependent FP8 WGMMAs.

In the context of attention, these layout restrictions entail certain modifications to the design of an FP8 algorithm, which we describe in \cref{sec:algofp8}.

\subsection{Standard Attention and Flash Attention} In line with the definitions established by \citet{dao2022flashattention}, we refer to \textbf{standard attention} as the approach where the GPU implementation generates and stores the intermediate matrices $\mathbf{S}$ and $\mathbf{P}$ directly into HBM. The principal innovation of \textsc{FlashAttention} involves utilizing a localized softmax reduction, thereby eliminating the need for costly intermediate memory accesses and consolidating the attention computation into a single kernel invocation. Specifically, this local softmax occurs within lines \ref{code-ws:softmax_start}-\ref{code-ws:softmax_end} of the consumer mainloop presented in \cref{alg:flash3_wgmma_ws_only}, in conjunction with rescaling blocks of $\mathbf{O}$. A concise derivation confirming that this procedure accurately computes $\mathbf{O}$ is provided in \cite[\S 2.3.1]{dao2023flashattention2}.

\section{FlashAttention-3: Algorithm}
\label{sec:algo}

In this section, we present the \textsc{FlashAttention-3} algorithm. Our discussion primarily addresses the forward pass, while details regarding the backward pass are provided in~\cref{sec:algo_ws_bwd}. Initially, we outline the method for combining warp-specialization with a circular SMEM buffer within the framework of the original \textsc{FlashAttention-2} algorithm. Subsequently, we elaborate on leveraging the asynchronous properties of WGMMA to establish a two-stage GEMM-softmax pipeline with overlapping operations. Lastly, we discuss the adaptations necessary for FP8, focusing on both layout compliance and achieving accuracy through block quantization and the use of incoherent processing techniques.

\subsection{Producer-Consumer asynchrony through warp-specialization and pingpong scheduling}
\label{sec:algo_ws}

\paragraph{Warp-specialization}
Similar to the approach in \textsc{FlashAttention-2}, the forward pass in \textsc{FlashAttention-3} can be easily parallelized across the batch dimension, head count, and query sequence length. Therefore, we focus on providing a CTA-level abstraction, wherein each CTA processes a specific tile $\mathbf{Q}_i$ of the query matrix and computes the matching tile $\mathbf{O}_i$ in the output. For clarity, our exposition first outlines the warp-specialization method that uses a circular SMEM buffer and omits any additional overlapping between GEMM and softmax operations. Denote by $d$ the dimensionality of each head and by $N$ the full sequence length. Given a selected query block size $B_r$, we partition $\mathbf{Q}$ into $T_r = \lceil \frac{N}{B_r} \rceil$ query tiles, namely $\mathbf{Q}_1, \ldots, \mathbf{Q}_{T_r}$.

\begin{algorithm}[H]
    \caption{\small\label{alg:flash3_wgmma_ws_only}\textsc{FlashAttention-3} forward computation \textbf{excluding} intra-consumer overlap -- CTA-centric structure}
    \begin{algorithmic}[1]
\REQUIRE Matrices $\mathbf{Q}_i \in \mathbb{R}^{B_r \times d}$ and $\mathbf{K}, \mathbf{V} \in \mathbb{R}^{N \times d}$ reside in HBM, with the key block dimension $B_c$ and total key blocks $T_c = \lceil \frac{N}{B_c} \rceil$.
\STATE Set up the pipeline controller for multi-stage barrier handling and a circular SMEM buffer of depth $s$.
\IF {current warpgroup is designated as producer}
\STATE Free a predefined set of registers.
\STATE Initiate transfer of $\mathbf{Q}_i$ from HBM to SMEM.
\STATE On completion, signal to consumers that $\mathbf{Q}_i$ is available in SMEM.
\FOR{$0 \le j < T_c$}
    \STATE Wait until consumers finish accessing the $(j\,\%\,s)$ slot in the buffer.
    \STATE Fetch $\mathbf{K}_j, \mathbf{V}_j$ from HBM into the $(j\,\%\,s)$ position in SMEM.
    \STATE Once loaded, notify consumers of the availability of $\mathbf{K}_j, \mathbf{V}_j$.
\ENDFOR
\ELSE \STATE Allocate a prescribed number of registers, adjusted to the count of consumer warps.
\STATE Initialize, in on-chip storage, $\mathbf{O}_i = (0) \in \mathbb{R}^{B_r \times d}$, and $\ell_i, m_i = (0), (-\infty) \in \mathbb{R}^{B_r}$.
\STATE Await the transfer of $\mathbf{Q}_i$ into SMEM.
\FOR{$0 \le j < T_c$}
\STATE Block until $\mathbf{K}_j$ becomes available in SMEM.
\STATE Perform SS-GEMM: $\mathbf{S}_i^{(j)} = \mathbf{Q}_i \mathbf{K}_j^T$. Commit and synchronize. \label{alg:ws_only_gemm1}
\STATE Save $m_i^{\mathrm{old}} = m_i$ and update $m_i = \mathrm{max}(m_i^{\mathrm{old}}, \mathrm{rowmax}(\mathbf{S}_i^{(j)}))$. \label{code-ws:softmax_start}
\STATE Compute $\widetilde{\mathbf{P}}_i^{(j)} = \mathrm{exp}(\mathbf{S}_i^{(j)} - m_i)$, and update $\ell_i = \mathrm{exp}(m_i^{\mathrm{old}} - m_i) \ell_i + \mathrm{rowsum}(\widetilde{\mathbf{P}}_i^{(j)})$. \label{code-ws:softmax_end}
\STATE Wait for $\mathbf{V}_j$ to finish loading into SMEM.
\STATE Conduct RS-GEMM: update $\mathbf{O}_i = \mathrm{diag}(\mathrm{exp}(m_i^{\mathrm{old}} - m_i))^{-1} \mathbf{O}_i + \widetilde{\mathbf{P}}_i^{(j)} \mathbf{V}_j$. Commit, then synchronize. \label{alg:ws_only_gemm2}
\STATE Mark the $(j\,\%\,s)$ buffer segment as available for producer use.
\ENDFOR
\STATE Apply the final normalization: $\mathbf{O}_i = \mathrm{diag}(\ell_i)^{-1} \mathbf{O}_i$ and compute $L_i = m_i + \log(\ell_i)$.
\STATE Store $\mathbf{O}_i$ and $L_i$ back into HBM as the $i$th blocks of $\mathbf{O}$ and $L$ respectively.
\ENDIF
\end{algorithmic}
\end{algorithm}

In our deployment of \cref{alg:flash3_wgmma_ws_only} on Hopper, we make use of \verb|setmaxnreg| to manage (de)allocation processes, employ TMA to load $\mathbf{Q}_i$ and $\{ \mathbf{K}_j, \mathbf{V}_j \}_{0 \leq j < T_c}$, and leverage WGMMA for performing the GEMMs within the consumer mainloop. Here, the prefixes SS or RS specify whether the first GEMM operand resides in shared memory or the register file, respectively. When analyzing the control flow of \cref{alg:flash3_wgmma_ws_only}, it is important to observe that TMA load instructions are asynchronous; consequently, dispatching a TMA load does not block the initiation of others. Furthermore, during the producer mainloop, the first $s$ iterations are devoted to filling the buffer, and no wait instructions are inserted during these initial cycles.

\paragraph{Pingpong scheduling} The inherent asynchronicity between WGMMA and TMA, combined with warp-specialization, provides an avenue for hiding the softmax computation latency of one warpgroup behind the GEMM execution of another. To illustrate the need for this strategy, note that on contemporary hardware accelerators, non-matmul operations exhibit considerably lower throughput relative to matmul operations. For instance, the H100 SXM5 GPU can perform matmul at 989 TFLOPS for FP16 precision, whereas its throughput for special functions such as exponentials\footnote{According to the CUDA programming guide, each streaming multiprocessor (SM) can process 16 special function operations per clock cycle. With 132 SMs running at a 1830 MHz clock speed, this yields approximately 3.9 TFLOPS for such functions.}—essential for softmax—reaches only 3.9 TFLOPS. In the context of an FP16 attention forward pass with a head dimension of 128, the volume of matmul FLOPs exceeds the number of exponential operations by a factor of 512, while the throughput for exponentials is 256 times lower; consequently, the time required for exponential computations may amount to 50\% of the time needed for matmuls. This imbalance becomes even more pronounced in FP8 computations, where matmul throughput is doubled yet the throughput for exponentials remains unchanged.

Given that the computation of the exponential function is handled by an independent hardware block—the multi-function unit—we aim to overlap the exponential computation with the matrix multiplication operations being executed by the Tensor Cores. To accomplish this, we introduce synchronization barriers (\texttt{bar.sync} instructions) to ensure that the GEMM operations (with GEMM1 corresponding to $\mathbf{P} \mathbf{V}$ within the current iteration, and GEMM0 to $\mathbf{Q} \mathbf{K}^\top$ for the subsequent iteration) in warpgroup 1 are launched prior to those in warpgroup 2. Consequently, while warpgroup 2 is engaged with its GEMM computations, the softmax calculation for warpgroup 1 proceeds in parallel. The roles then alternate: warpgroup 1 computes the GEMMs as warpgroup 2 performs softmax, thus establishing a “pingpong” scheduling mechanism as depicted in~\cref{fig:pingpong_scheduling}. Although, in real systems, this interleaving is somewhat less idealized than represented in the illustration, we consistently observe a performance benefit from this approach (for example, boosting measured throughput from 570 TFLOPS to 620–640 TFLOPS on FP16 forward pass with a head size of 128 and a sequence length of 8192).

\paragraph{Attention variants} When handling multi-query attention~\citep{shazeer2019fast} and grouped query attention~\citep{ainslie2023gqa}, we utilize the strategy outlined in \textsc{FlashAttention-2}, modifying tensor indexing such that $\mathbf{K}$ and $\mathbf{V}$ are not redundantly stored in HBM.

\subsection{Intra-warpgroup overlapping GEMMs and softmax}

Even within one warpgroup, we can overlap some instructions in the softmax with some instructions in the GEMMs. We describe one technique to do so.

Within the attention mechanism, the main loop presents sequential dependencies that hinder the possibility of parallel execution within a single iteration. As an illustration, (local) softmax computations (see lines \ref{code-ws:softmax_start} to \ref{code-ws:softmax_end}) depend on the values of $\mathbf{S}_i^{(j)}$ produced by the initial GEMM, while the subsequent GEMM operation uses $\widetilde{\mathbf{P}}_i^{(j)}$—the output of the softmax—as input. Consequently, the wait statements at lines \ref{alg:ws_only_gemm1} and \ref{alg:ws_only_gemm2} in \cref{alg:flash3_wgmma_ws_only} enforce a serialized flow between softmax and GEMM calculations. Nevertheless, these sequential dependencies may be alleviated by employing pipelining across loop iterations, facilitated by registering extra intermediate buffers. Building upon this insight, we introduce the following two-stage\footnote{Observe that, although the degree of overlap is constrained by the number $s$ of circular SMEM buffer stages, it does not have to match it precisely.} pipeline algorithm, which integrates both GEMM and softmax steps:

\begin{algorithm}[H]
    \caption{\small\label{alg:flash3_wgmma}\textsc{FlashAttention-3} consumer warpgroup forward pass}
    \begin{algorithmic}[1]
      \REQUIRE  Given: matrices $\mathbf{Q}_i \in \mathbb{R}^{B_r \times d}$, $\mathbf{K}, \mathbf{V} \in \mathbb{R}^{N \times d}$ stored in HBM; key block size $B_c$ is defined and the total number of key blocks is $T_c = \lceil \frac{N}{B_c} \rceil$.
      \STATE Adjust register allocation according to the quantity of consumer warps utilized.
      \STATE Allocate and initialize on-chip variables as follows: set $\mathbf{O}_i = (0) \in \mathbb{R}^{B_r \times d}$; assign initial values $\ell_i, m_i = (0), (-\infty) \in \mathbb{R}^{B_r}$.
      \STATE Hold until both $\mathbf{Q}_i$ and $\mathbf{K}_0$ are present in shared memory.
      \STATE Employ WGMMA to compute $\mathbf{S}_{\mathrm{cur}} = \mathbf{Q}_i \mathbf{K}_0^T$. Commit the operation and synchronize for completion.
      \STATE Release the initial stage of the buffer allocated for $\mathbf{K}$.
      \STATE Calculate $m_{i}$, $\tilde{\mathbf{P}}_{\mathrm{cur}}$, and $\ell_{i}$ based on the values from $\mathbf{S}_{\mathrm{cur}}$, then appropriately rescale $\mathbf{O}_i$.
      \FOR{$1 \le j < T_c - 1$}
        \STATE Wait until $\mathbf{K}_j$ is loaded to shared memory.
        \label{code:mainloop_start}
        \STATE Perform WGMMA to determine $\mathbf{S}_{\mathrm{next}} = \mathbf{Q}_i \mathbf{K}_{j}^T$; commit without waiting for its completion. \label{code:first_wgmma}
        \STATE Pause until $\mathbf{V}_{j-1}$ becomes available in shared memory.
        \STATE Update $\mathbf{O}_{i}$ by accumulating $\tilde{\mathbf{P}}_{\mathrm{cur}} \mathbf{V}_{j-1}$ using WGMMA, commit and continue without waiting. \label{code:second_wgmma}
        \STATE Synchronize for completion of WGMMA operation involving $\mathbf{Q}_i \mathbf{K}_{j}^T$.
        \STATE Derive $m_{i}$, $\tilde{\mathbf{P}}_{\mathrm{next}}$, and $\ell_{i}$ from $\mathbf{S}_{\mathrm{next}}$. \label{code:softmax}
        \STATE Wait for WGMMA output of $\tilde{\mathbf{P}}_{\mathrm{cur}} \mathbf{V}_{j-1}$ to finalize, then update and rescale $\mathbf{O}_i$ accordingly.
        \STATE Free the $(j\,\%\,s)$th stage (for $\mathbf{K}$ buffer) and the $(j-1\,\%\,s)$th stage (for $\mathbf{V}$ buffer), respectively.
        \STATE Overwrite $\mathbf{S}_{\mathrm{cur}}$ with data from $\mathbf{S}_{\mathrm{next}}$.
        \label{code:mainloop_end}
      \ENDFOR \STATE Wait until $\mathbf{V}_{T_c - 1}$ has been loaded into shared memory.
      \STATE Conduct the computation $\mathbf{O}_{i} = \mathbf{O}_{i} + \tilde{\mathbf{P}}_{\mathrm{last}} \mathbf{V}_{T_c - 1}$ via WGMMA. Commit this operation and ensure its completion.

\cref{alg:flash3_wgmma} serves to substitute the consumer pathway in \cref{alg:flash3_wgmma_ws_only}, thereby completing the implementation of the \textsc{FlashAttention-3} method for FP16 computations. In a broad sense, we use WGMMA to refer generally to asynchronous GEMM. Throughout the principal loop (lines \ref{code:mainloop_start} through \ref{code:mainloop_end}), the execution of the second WGMMA kernel in iteration $j$ (line \ref{code:second_wgmma}) is performed in parallel with the softmax computations of the succeeding iteration, $j+1$ (line \ref{code:softmax}).

Although the pipelined scheme depicted earlier promises significant theoretical speedup, its real-world effectiveness depends on several implementation details:
\paragraph{Compiler reordering} The pseudocode assumes a specific ordering of operations, but in reality, the compiler (NVCC) frequently reorders instructions for performance reasons. This optimization may interfere with the intended alternation between WGMMA and other operations, potentially undermining the expected pipeline efficiency or altering behavior. Inspection of the SASS output verifies that overlapping instruction regions are indeed generated as intended (refer to Section~\ref{sec:2-stage-sass}).
\paragraph{Register pressure} Achieving maximum efficiency requires limiting register spilling. The 2-stage pipeline design necessitates additional registers, as intermediate outputs and pipeline context must be stored—most notably, another copy of $\mathbf{S}_{\mathrm{next}}$. This incurs an extra per-threadblock register usage of $B_r \times B_c \times \text{sizeof}(\text{float})$. Heightened register demands could restrict feasible block sizes, since both optimizations tend to exhaust the available register file. In practice, the optimal balance between them must be determined through empirical performance measurements.
\paragraph{3-stage pipelining} Building on the aforementioned 2-stage method, we extend to a 3-stage pipeline where the second WGMMA computation is more tightly overlapped with the softmax stage. This design might further enhance Tensor Core occupancy, but increases register needs even more, complicating the choice between greater tiling granularity and deeper pipelines. The full design and performance evaluation of this 3-stage approach are presented in ~\cref{sec:3-stage}.

\subsection{Low-precision with FP8}
\label{sec:algofp8}

\textbf{Efficiency: layout transformations.} Executing the forward pass of \textsc{FlashAttention-3} using FP8 precision introduces layout compatibility challenges beyond those seen with FP16.

Typically, the input tensors $\mathbf{Q}$, $\mathbf{K}$, and $\mathbf{V}$ are arranged contiguously along the head dimension. However, to meet the k-major layout requirement imposed by FP8 WGMMA in the second GEMM, it becomes necessary for $\mathbf{V}$—specifically, the tiles of $\mathbf{V}$ stored in SMEM—to be contiguous along the sequence length dimension. Since TMA loads cannot alter the contiguous dimension, two primary approaches emerge: (1) performing a pre-processing transpose of $\mathbf{V}$ in GMEM or (2) conducting an in-kernel transpose of $\mathbf{V}$'s tiles after loading them into SMEM. To pursue option (1), this transpose can be (1a) fused directly into the epilogue of a prior operation (such as applying rotary embedding), or (1b) executed as a standalone pre-processing transpose kernel\footnote{An optimized transpose kernel can operate close to the theoretical bandwidth of the hardware~\citep{colfax_cutlass_transpose_2024}.} that swaps the strides of the sequence length and head dimensions. However, integrating (1a) within a generic library is non-trivial, while (1b) is generally inefficient for memory-bound scenarios like inference, due to high memory overhead.

Rather, in FP8 \textsc{FlashAttention-3}, we select approach (2). For performing in-kernel transpositions, we leverage the LDSM (\verb|ldmatrix|) and STSM (\verb|stmatrix|) operations, which enable a warp of threads to jointly move data between SMEM and RMEM with a block size of 128 bytes. \footnote{According to the PTX documentation, LDSM/STSM manipulate $8 \times 8$ tiles comprising 16-bit elements \cite[\S 9.7.13.4.15-16]{ptx}; by packing two 8-bit elements together, these instructions can be repurposed for FP8 formats. Nevertheless, the transpose functionality in LDSM/STSM does not support unpacking of 8-bit entries, requiring specific shuffling of register data between their use to achieve tile-level transposition, though we do not detail these operations here.} The LDSM and STSM primitives are highly efficient in their use of registers, making it feasible to perform these operations within the producer warpgroup, and additionally support layout transposition as part of the memory transfer process. Furthermore, after the initial step, it is possible to schedule the transposition of the subsequent $\mathbf{V}$ tile such that it is overlapped with the execution of the two WGMMAs pertaining to the previous $\mathbf{V}$ and current $\mathbf{K}$ tiles.

Second, we note that, unlike FP16, the memory organization for the FP32 accumulator in an FP8 WGMMA does not align with the arrangement expected for operand A when stored in registers. Illustrations of these two respective register layouts are shown in \cref{fig:rmem_accum} and \cref{fig:rmem_operand}, where each register entry corresponds to a per-thread value following the displayed order. Through the use of byte permute instructions, it is possible to convert the accumulator output from the initial WGMMA into an arrangement appropriate for input to a subsequent WGMMA, which is also consistent with the register order of the $\mathbf{V}$ tile as produced by an in-kernel transpose. Specifically, referring to \cref{fig:rmem_accum}, we systematically reorder the entries as
$$\{ \verb|d0 d1 d4 d5 d2 d3 d6 d7| \},$$
repeating this register permutation every 8 bytes. On the logical level, this procedure effectively reorders the columns of the $\mathbf{P}$ tile (for instance, columns $0189$ are mapped to the first four positions). To ensure WGMMA generates the correct output tile, we may similarly tailor the in-kernel transpose so that the rows of the $\mathbf{V}$ tile are permuted accordingly.\footnote{This extra flexibility, made possible by in-kernel transposition, removes the need for shuffle instructions to transfer register values among threads, which had been required as discussed in~\citep{colfax_fp8_flashattention_2024}.}

\textbf{Accuracy: block quantization and incoherent processing.} The FP8 (e4m3) representation allocates 3 bits to the mantissa and 4 bits to the exponent, leading to increased numerical error relative to FP16 or BF16 formats. Additionally, large-scale models tend to exhibit significant outlier values~\citep{dettmers2208llm, sun2024massive}, which complicates effective quantization due to their substantial deviations in magnitude. Standard practice involves per-tensor scaling~\citep{micikevicius2022fp8}, wherein a separate scaling factor is assigned to each tensor (for instance, one each for $\mathbf{Q}$, $\mathbf{K}$, and $\mathbf{V}$).
To improve the numerical precision of attention operations in FP8, we utilize two strategies:

\begin{enumerate}

\item \textbf{Block quantization}: In this approach, a single scalar is retained for each block; specifically, the tensors $\mathbf{Q}$, $\mathbf{K}$, and $\mathbf{V}$ are partitioned into segments of size $B_r \times d$ or $B_c \times d$, and quantized independently per block. This quantization step can be seamlessly combined with the operation preceding attention (such as rotary embedding), without incurring any additional latency, as the latter is constrained by memory bandwidth. Due to the inherently block-based design of the \textsc{FlashAttention-3} method, each block of $\mathbf{S}$ can be scaled appropriately to adjust for block-wise quantization, incurring no extra computational overhead.

\item \textbf{Incoherent processing}: To mitigate the effects of outliers, prior to quantization into FP8, both $\mathbf{Q}$ and $\mathbf{K}$ are pre-multiplied with a random orthogonal matrix $\mathbf{M}$. Given the property $\mathbf{M} \mathbf{M}^\top = I$ for orthogonal matrices, it follows that $(\mathbf{Q} \mathbf{M}) (\mathbf{K} \mathbf{M})^\top = \mathbf{Q} \mathbf{K}^\top$, ensuring that the application of $\mathbf{M}$ leaves the resulting attention computation unaffected. This transformation helps to redistribute outlier magnitudes, since each element of $\mathbf{Q} \mathbf{M}$ or $\mathbf{K} \mathbf{M}$ becomes a random linear combination of elements from $\mathbf{Q}$ or $\mathbf{K}$, thereby lessening quantization errors. Following the methodology of \citet{chee2024quip} and~\citet{tseng2024quip}, $\mathbf{M}$ is selected as the product of random diagonal $\pm 1$ matrices and a Hadamard matrix. This structure allows multiplication in $O(d \log d)$ time, as opposed to the typical $O(d^2)$, and this operation can likewise be incorporated into rotary embedding without incurring further computation cost.
\end{enumerate}

We validate that these two techniques reduces numerical error by up to 2.6$\times$ in \cref{sec:numerical_error}.

\section{Empirical Validation}
\label{sec:exp}

We use the primitives from CUTLASS~\citep{Thakkar_CUTLASS_2023} such as WGMMA and TMA abstractions to implement \textsc{FlashAttention-3} and evaluate its efficiency and accuracy.

\begin{itemize}

\item \textbf{Benchmarking attention.} We evaluate the execution time of \textsc{FlashAttention-3} over a range of sequence lengths and compare its performance against the baseline PyTorch implementation, \textsc{FlashAttention-2}, \textsc{FlashAttention-2} implemented with Triton (which leverages H100-specific instructions), as well as a vendor-provided \textsc{FlashAttention-2} version optimized for H100 GPUs using cuDNN. Our experiments demonstrate that \textsc{FlashAttention-3} achieves speedups of up to 2.0$\times$ relative to \textsc{FlashAttention-2}, and is up to 1.5$\times$ faster than the Triton-optimized version. Notably, \textsc{FlashAttention-3} attains 740 TFLOPs/s, representing 75\% of the H100 GPU's peak theoretical throughput.

\item \textbf{Ablation study.} Our results show that enhancements such as warp-specialization and GEMM-softmax pipelining are key factors driving the improved efficiency observed in \textsc{FlashAttention-3}.

\item \textbf{Accuracy of FP8 attention.} We demonstrate that applying block quantization along with incoherent processing leads to a 2.6$\times$ reduction in numerical errors for FP8 \textsc{FlashAttention-3}.
\end{itemize}

\subsection{Benchmarking Attention}
\label{subsec:benchmark_attn}

We evaluate the execution times of various attention mechanisms on an H100 80GB SXM5 GPU across different configurations (presence or absence of causal mask, head dimensions of 64 and 128) using FP16 precision. Our findings, presented in~\cref{fig:benchmark_attn_fwd} and~\cref{fig:benchmark_attn_bwd}, indicate that \textsc{FlashAttention-3} outperforms \textsc{FlashAttention-2} by approximately 1.5-2.0$\times$ in the forward direction and 1.5-1.75$\times$ in the backward direction. When set against a conventional attention baseline, \textsc{FlashAttention-3} achieves speedups ranging from 3$\times$ to 16$\times$. Notably, for sequence lengths at or above 1k, \textsc{FlashAttention-3} is even able to exceed the throughput of the proprietary vendor library (cuDNN), which is tailored for the H100 GPUs.

\paragraph{Benchmark settings:} The sequence lengths considered are 512, 1k, up to 16k, while the batch size is chosen to ensure a constant total token count of 16k. The hidden dimension is fixed at 2048, and the head dimension is selected from 64, 128, or 256, which corresponds to configurations of 32, 16, or 8 heads, respectively. For the computation of forward pass FLOPs, we employ:

\begin{equation*}
  4 \cdot \text{seqlen}^2 \cdot \text{head dimension} \cdot \text{number of heads}.
\end{equation*}

Applying causal masking reduces the number of computed entries by approximately half, so we divide the total by 2 to reflect this. For the backward pass, the required FLOPs are estimated by multiplying the forward pass FLOPs by 2.5, considering that while the forward computation involves 2 matrix multiplications, the backward step comprises 5 matmuls as a result of recomputation.

We additionally assess the forward pass runtime in FP8 precision using the same setup. The outcomes for \texttt{headdim} 256 are illustrated in ~\cref{fig:benchmark_attn_fp8}, with the comprehensive results available in \cref{sec:benchmark_attn_fp8_full}.

\subsection{Ablation Study: 2-Stage Pipelining Experiments}
\label{sec:2-stage-pipelining-experiments}

We conduct an ablation study on both the 2-stage WGMMA-softmax pipeline and warp-specialization for non-causal FP16 \textsc{FlashAttention-3}, utilizing fixed parameters $\{\text{batch}, \text{seqlen}, \text{nheads}, \text{hdim}\} = \{ 4, 8448, 16, 128\}$. As shown in~\cref{table:ablation_pipelining}, the results validate that our enhancements—incorporating asynchrony through warp-specialization and allowing GEMM and softmax operations to overlap—yield a substantial performance boost, improving throughput from 570 to 661 TFLOPs.

\subsection{Numerical Error Validation}
\label{sec:numerical_error}

Given the growing attention to the numerical error associated with \textsc{FlashAttention}~\citep{golden2024flash}, we evaluate \textsc{FlashAttention-2}, \textsc{FlashAttention-3}, and a conventional attention implementation, comparing each to a reference implementation in FP64. To mimic the presence of outlier features and activations typical in LLMs~\citep{dettmers2208llm, sun2024massive}, we sample the elements of $\mathbf{Q}, \mathbf{K}, \mathbf{V}$ using the following distribution:

\begin{equation*}
  \mathcal{N}(0, 1) + \mathcal{N}(0, 100) \cdot \mathrm{Bernoulli}(0.001).
\end{equation*}

Specifically, each element follows a normal distribution with a mean of zero and a standard deviation of 1; however, in 0.1\% of entries, an independent normal component with standard deviation 10 is added. The root mean squared error (RMSE) results are reported in~\cref{table:numerical_error}. In FP16 precision, both \textsc{FlashAttention-2} and \textsc{FlashAttention-3} reduce RMSE by a factor of 1.7 relative to the conventional implementation, as they retain intermediate (softmax) computations in FP32. For FP8, the baseline employs per-tensor scaling, uses FP32 for the matmul accumulator, and stores the softmax intermediates in FP16. Owing to the adoption of block quantization and incoherent processing, \textsc{FlashAttention-3} in FP8 offers RMSE that is 2.6$\times$ lower than this baseline.

\section{Dicussion, Limitations, Conclusion}
\label{sec:discussion}

Through \textsc{FlashAttention-3}, we illustrate that adopting novel programming paradigms and leveraging hardware innovations like asynchrony and low-precision computation substantially enhances both the performance and reliability of attention mechanisms. Our approach achieves a 1.5–2.0$\times$ speedup over \textsc{FlashAttention-2}, and reduces FP8-related numerical inaccuracies by a factor of 2.6 compared to conventional per-tensor quantization. However, several areas remain for future improvement, such as fine-tuning performance for LLM inference, incorporating a persistent kernel architecture within the FP8 kernel,\footnote{In our experimental setup, the FP16 version of \textsc{FlashAttention-3} utilizes a persistent kernel and an advanced load balancing method, whereas the FP8 implementation does not, which partly accounts for the inferior FP8 \textsc{FlashAttention-3} performance at shorter sequence lengths and when using causal masking, relative to FP8 cuDNN kernels.} and gaining a deeper understanding of low-precision attention in the context of large-scale model training. While this study focuses on Hopper GPUs, we anticipate that our methods will be applicable to a broad array of hardware accelerators. Ultimately, we believe that more efficient and precise attention primitives can pave the way for progress in long-context problem domains.

\end{document}